\section{Программа}

\subsection{База знаний}
\begin{listing}[Исходный код базы знаний]{lst:t1-prog}{Prolog}
len(List, Length) :- 
    len_rec(List, 0, Length).

len_rec([], Acc, Acc) :- !.
len_rec([_|Tail], Acc, Length) :- 
    NewAcc is Acc + 1,
    len_rec(Tail, NewAcc, Length).


sum(List, Sum) :- 
    sum_rec(List, 0, Sum).

sum_rec([], Acc, Acc) :- !.
sum_rec([Head|Tail], Acc, Sum) :- 
    NewAcc is Acc + Head,
    sum_rec(Tail, NewAcc, Sum).


sum_odd_pos(List, Sum) :- 
    sum_odd_pos_rec(List, 0, Sum).

sum_odd_pos_rec([], Acc, Acc) :- !.
sum_odd_pos_rec([_], Acc, Acc) :- !.
sum_odd_pos_rec([_, Odd | Tail], Acc, Sum) :- 
    NewAcc is Acc + Odd,
    sum_odd_pos_rec(Tail, NewAcc, Sum).


filter_gt(List, V, Result) :- 
    filter_gt_rec(List, V, [], Result).

filter_gt_rec([], _, Acc, Acc) :- !.
filter_gt_rec([Head|Tail], V, Acc, Result) :- 
    Head > V, !, 
    filter_gt_rec(Tail, V, [Head|Acc], Result).
filter_gt_rec([_|Tail], V, Acc, Result) :- 
    filter_gt_rec(Tail, V, Acc, Result).


remove(List, Element, Result) :- 
    remove_rec(List, Element, [], Result).

remove_rec([], _, Acc, Acc) :- !.
remove_rec([Element|Tail], Element, Acc, Result) :- !,
    remove_rec(Tail, Element, Acc, Result).
remove_rec([Head|Tail], Element, Acc, Result) :- 
    remove_rec(Tail, Element, [Head|Acc], Result).


merge(List1, List2, Result) :- 
    prepend(List1, List2, Result).

prepend([], AccList, AccList) :- !.
prepend([Head|Tail], AccList, Result) :- 
    prepend(Tail, [Head|AccList], Result).
\end{listing}

\subsection{Вопросы задания}

Для получения длины списка необходимо задать вопрос: 
существует ли такой Результат при конкретизированном Списке, для которого 
удовлетворяется отношение \code{len}.

\begin{listing}[Нахождение длины списка]{lst:q1}{Prolog}
?- len([a, b, c, d], L).
\end{listing}

Для получения суммы элементов числового списка необходимо задать 
вопрос: существует ли такой Результат при конкретизированном числовом 
Списке, для которого удовлетворяется отношение \code{sum}.

\begin{listing}[Нахождение суммы элементов списка]{lst:q2}{Prolog}
?- sum([1, 2, 3, 4], S).
\end{listing}

Для получения суммы элементов числового списка, стоящих на 
нечетных позициях (нумерация с 0), необходимо задать вопрос: 
существует ли такой Результат при конкретизированном 
числовом Списке, для которого удовлетворяется 
отношение \code{sum\_odd\_pos}.

\begin{listing}[Сумма элементов на нечетных позициях]{lst:q3}{Prolog}
?- sum_odd_pos([10, 20, 30, 40, 50], S).
\end{listing}

Для формирования списка из элементов числового списка, больших 
заданного значения, необходимо задать вопрос: существует ли такой 
результирующий Список при конкретизированных исходном числовом Списке 
и пороговом Значении, для которых удовлетворяется отношение \code{filter\_gt}.

\begin{listing}[Фильтрация списка по пороговому значению]{lst:q4}{Prolog}
?- filter_gt([2, 8, 3, 10, 5], 5, Res).
\end{listing}

Для удаления заданного элемента из списка необходимо задать вопрос: 
существует ли такой результирующий Список при конкретизированных 
исходном Списке и удаляемом Элементе, для которых удовлетворяется 
отношение \code{remove}.

\begin{listing}[Удаление заданного элемента из списка]{lst:q5}{Prolog}
?- remove([1, 3, 2, 3, 4], 3, Res).
\end{listing}

Для объединения двух списков в один результирующий необходимо 
задать вопрос: существует ли такой результирующий Список при 
конкретизированных первом Списке и втором Списке, для которых 
удовлетворяется отношение \code{merge}.

\begin{listing}[Объединение двух списков]{lst:q6}{Prolog}
?- merge([1, 2], [3, 4], Res).
\end{listing}